\documentclass[12pt]{article}
\usepackage[utf8]{inputenc}
\usepackage[T1]{fontenc}
\usepackage{hyperref}
\title{Solomon Islands Sign Language Language Resources for K-12 Curriculum Development}
\author{Compiled by Claude Sonnet 5.5 (claude-sonnet-5-5) for Andr\'as Kornai}
\date{}
\begin{document}
\maketitle

\section{Introduction}
Solomon Islands Sign Language (ISO 639-3 szs) is a deaf community sign language of the Solomon Islands. This report is based on a single round of web searches. Very little was found, and nothing that is a curriculum or a published description of the language.

\section{Primary Educational Resources}
No textbook, primer or school curriculum in Solomon Islands Sign Language was found. Press reports describe an education project for deaf youth, Iumi Tugeda Helpem Solomon Islands Def Comuniti, launched at the San Isidro Care Centre. It aims to train sign language trainers and to give teachers professional development on deaf language and inclusive education. See \url{https://www.sibconline.com.sb/partners-launch-ground-breaking-project-supporting-education-for-deaf-youth}.

A project by Australian Catholic University and the Archdiocese of Honiara, funded through the Education Sector Support Program, trains deaf and hearing people in Auslan and sends them to remote communities. It offers a Certificate of Teaching and Learning recognised by the Ministry of Education. See \url{https://acu.edu.au/about-acu/news/2022/july/new-project-will-send-sign-language-teachers-into-remote-villages-of-solomon-islands}. The project teaches Auslan, which is a different language from Solomon Islands Sign Language, so it does not by itself supply local-language materials.

\section{Linguistic and Reference Materials}
No dictionary, grammar or corpus of Solomon Islands Sign Language was found in these searches.

\section{Religious and Community Resources}
The only educational facility for deaf youth reported is the San Isidro Care Centre in Honiara, run by the Catholic Education Authority. See \url{https://thesoutherncross.org.au/news/2022/07/15/support-for-deaf-youth-in-solomon-islands/}.

\section{Digital and Audio Resources}
No video or digital resources in the language were found.

\section{Curriculum Development Context}
The reports say that no research establishes how many deaf people there are in the country. The Ministry of Education and Human Resources Development reported recognising 16 sign language trainers. They taught basic sign language to 498 school children and community members, and the Ministry distributed sign language books and posters. A curriculum would have to be built almost from scratch, together with the deaf community. The 2022 reports may be out of date.
\end{document}